\chapter*{Abstract}

Modern relational databases increasingly suffer from data quality problems
including inconsistencies, duplicates, missing values, and constraint violations
arising from heterogeneous data sources and evolving business rules. Existing
solutions, ranging from interactive visual tools and SQL-based rule engines to
enterprise data quality platforms and DSL-based pipeline frameworks, either embed
validation logic in graphical interfaces or proprietary configurations, or focus
on data transformation pipelines rather than on an explicit, reusable, and
auditable specification of data quality constraints. This paper presents DaQuVa
(Data Quality Validator), a novel domain-specific language designed to address
these gaps by providing a declarative, safety-first approach to data quality
enforcement in SQL databases. DaQuVa introduces a formally defined grammar, a
sequential execution model, and a constrained expression language that makes
programs straightforward to review and audit. The language provides core
constructs for column-level scanning via pluggable validation tools,
threshold-based filtering, local duplicate detection with optional AI-powered
similarity scoring, controlled database writes through an explicit safety
mechanism (\texttt{allowDanger}), result export to multiple sinks, and modular
reusability through parameterized functions. A Python-based implementation
comprising a regex-driven lexer, a recursive-descent parser, and an execution
engine is described. Experimental evaluation demonstrates that DaQuVa can detect
and flag erroneous records and near-duplicate entries across realistic datasets
with concise, human-readable scripts, while its safety constraints prevent
accidental data loss. Unlike general-purpose DSLs or tool-driven approaches,
DaQuVa treats data quality rules as a formal, portable contract between technical
teams and domain experts, contributing a practical foundation for reproducible
and auditable data quality management.
